\documentclass[a4paper, 12pt]{report}		% general format
\usepackage{multicol}
%%%% Charset
\usepackage{cmap}							% make PDF files searchable and copyable
\usepackage{bm}
\usepackage{pdfpages}
\usepackage[utf8x]{inputenc} 				% accept different input encodings
\usepackage[english,russian]{babel}   %% загружает пакет многоязыковой вёрстки
%\usepackage{fontspec}      %% подготавливает загрузку шрифтов Open Type, True Type и др.
%\defaultfontfeatures{Ligatures={TeX},Renderer=Basic}  %% свойства шрифтов по умолчанию
%\setmainfont[Ligatures={TeX,Historic}]{Roboto-Light} %% задаёт основной шрифт документа
%\setsansfont{Roboto-Light}  
\usepackage{float}
%%%% Graphics
%\usepackage[dvipsnames]{xcolor}			% driver-independent color extensions
\usepackage{graphicx}						% enhanced support for graphics
\usepackage{wrapfig}						% produces figures which text can flow around

%%%% Math
\usepackage{amsmath}						% American Mathematical Society (AMS) math facilities
\usepackage{amsfonts}						% fonts from the AMS
\usepackage{amssymb}						% additional math symbols

%%%% Typograpy (don't forget about cm-super)
\usepackage{microtype}						% subliminal refinements towards typographical perfection
\linespread{1.0}							% line spacing
\usepackage[mag=1000, left=2.0cm, right=1.5cm, top=2cm, bottom=2cm, headsep=0.7cm, footskip=1cm]{geometry}
\setlength{\parindent}{0pt}					% we don't want any paragraph indentation
\usepackage{parskip}						% some distance between paragraphs

%%%% Tables
\usepackage{tabularx}						% tables with variable width columns
\usepackage{multirow}						% for tabularx
\usepackage{hhline}							% for tabularx
\usepackage{tabu}
\usepackage{longtable}

%%%% Graph
\usepackage{tikz}							% package for creating graphics programmatically
\usetikzlibrary{arrows}						% edges for tikz

%%%% Other
\usepackage{url}							% verbatim with URL-sensitive line breaks
\usepackage{fancyvrb}						% sophisticated verbatim text (with box)

\usepackage{fancyhdr}
\usepackage{latexsym}
\usepackage{booktabs}
\usepackage{array}

\usepackage{listings}
\usepackage{caption}
\DeclareCaptionFont{white}{\color{white}}
\DeclareCaptionFormat{listing}{\colorbox{gray}{\parbox{\dimexpr\textwidth-1.72\fboxsep\relax}{#1#2#3}}}
\captionsetup[lstlisting]{format=listing,labelfont=white,textfont=white,margin=0pt}
\lstset{language=C,
	basicstyle=\footnotesize,
	keepspaces=true,
	tabsize=4,               
	frame=single,                           % Single frame around code
	rulecolor=\color{black},
	captionpos=b,
	showstringspaces=false,	
	abovecaptionskip=-0.9pt,
	xleftmargin=3.4pt,
	xrightmargin=2.6pt,
	breaklines=true,
	postbreak=\raisebox{0ex}[0ex][0ex]{\ensuremath{\color{black}\hookrightarrow\space}},
	xleftmargin=3.2pt,
	literate={а}{{\selectfont\char224}}1
	{~}{{\textasciitilde}}1
	{б}{{\selectfont\char225}}1
	{в}{{\selectfont\char226}}1
	{г}{{\selectfont\char227}}1
	{д}{{\selectfont\char228}}1
	{е}{{\selectfont\char229}}1
	{ё}{{\"e}}1
	{ж}{{\selectfont\char230}}1
	{з}{{\selectfont\char231}}1
	{и}{{\selectfont\char232}}1
	{й}{{\selectfont\char233}}1
	{к}{{\selectfont\char234}}1
	{л}{{\selectfont\char235}}1
	{м}{{\selectfont\char236}}1
	{н}{{\selectfont\char237}}1
	{о}{{\selectfont\char238}}1
	{п}{{\selectfont\char239}}1
	{р}{{\selectfont\char240}}1
	{с}{{\selectfont\char241}}1
	{т}{{\selectfont\char242}}1
	{у}{{\selectfont\char243}}1
	{ф}{{\selectfont\char244}}1
	{х}{{\selectfont\char245}}1
	{ц}{{\selectfont\char246}}1
	{ч}{{\selectfont\char247}}1
	{ш}{{\selectfont\char248}}1
	{щ}{{\selectfont\char249}}1
	{ъ}{{\selectfont\char250}}1
	{ы}{{\selectfont\char251}}1
	{ь}{{\selectfont\char252}}1
	{э}{{\selectfont\char253}}1
	{ю}{{\selectfont\char254}}1
	{я}{{\selectfont\char255}}1
	{А}{{\selectfont\char192}}1
	{Б}{{\selectfont\char193}}1
	{В}{{\selectfont\char194}}1
	{Г}{{\selectfont\char195}}1
	{Д}{{\selectfont\char196}}1
	{Е}{{\selectfont\char197}}1
	{Ё}{{\"E}}1
	{Ж}{{\selectfont\char198}}1
	{З}{{\selectfont\char199}}1
	{И}{{\selectfont\char200}}1
	{Й}{{\selectfont\char201}}1
	{К}{{\selectfont\char202}}1
	{Л}{{\selectfont\char203}}1
	{М}{{\selectfont\char204}}1
	{Н}{{\selectfont\char205}}1
	{О}{{\selectfont\char206}}1
	{П}{{\selectfont\char207}}1
	{Р}{{\selectfont\char208}}1
	{С}{{\selectfont\char209}}1
	{Т}{{\selectfont\char210}}1
	{У}{{\selectfont\char211}}1
	{Ф}{{\selectfont\char212}}1
	{Х}{{\selectfont\char213}}1
	{Ц}{{\selectfont\char214}}1
	{Ч}{{\selectfont\char215}}1
	{Ш}{{\selectfont\char216}}1
	{Щ}{{\selectfont\char217}}1
	{Ъ}{{\selectfont\char218}}1
	{Ы}{{\selectfont\char219}}1
	{Ь}{{\selectfont\char220}}1
	{Э}{{\selectfont\char221}}1
	{Ю}{{\selectfont\char222}}1
	{Я}{{\selectfont\char223}}1,
	extendedchars=true
}

%галочка
\usepackage{amssymb}% http://ctan.org/pkg/amssymb
\usepackage{pifont}% http://ctan.org/pkg/pifont
\newcommand{\cmark}{\ding{52}}%
\newcommand{\xmark}{\ding{56}}
%------------------------------------------------------------------------------
\renewcommand{\labelenumii}{\theenumii}
\renewcommand{\theenumii}{\theenumi.\arabic{enumii}.}
\addto\captionsrussian{\def\refname{Список использованных источников}}
\begin{document}
\begin{titlepage}
\thispagestyle{empty}

\begin{center}
Санкт-Петербургский политехнический университет Петра Великого\\
Институт Информационных Технологий и Управления\\*
Кафедра компьютерных систем и программных технологий\\*
\hrulefill
\end{center}

\vspace{15em}

\begin{center}
\textsc{\textbf{Курсовая работа}}
\vspace{1em}

Дисциплина: \textbf{Методы оптимизации}
\vspace{2em}

Тема: \textbf{Формулировка и решение задачи выбора оптимального решения с использованием различных математических моделей}
\end{center}

\vspace{16em}

\begin{flushleft}
Выполнил студент гр. 53501/3 \hrulefill С.А. Мартынов \\
\vspace{1.5em}
Руководитель, к.т.н.,доц. \hrulefill А.Г. Сиднев\\
\end{flushleft}

\vspace{\fill}

\begin{center}
Санкт-Петербург \\
2015
\end{center}

\end{titlepage}
\setcounter{page}{2}
%\includepdf[pages=-,pagecommand={},width=\textwidth]{task.pdf}
\tableofcontents
\clearpage

%------------------------------------------------------------------------------
%\input{intro}

\addcontentsline{toc}{chapter}{Введение}
\chapter*{Введение}

В данной курсовой работе рассматривается моделирование корпоративной компьютерной сети юридической фирмы средствами Cisco Packet Tracer \cite{cite-cis}. Работа разделена на несколько этапов:

\begin{itemize}
	\item построение корпоративной компьютерной сети в Cisco Packet Tracer \cite{cite-cis};
	\item настройка корпоративных сетевых сервисов;
	\item тестирование работоспособности корпоративных сетевых сервисов.
\end{itemize}

%%%%%%%%%%%%%%%%%%%%%%%%%%%%%%%%%%%%%%%%%
% Journal Article
% LaTeX Template
% Version 1.4 (15/5/16)
%
% This template has been downloaded from:
% http://www.LaTeXTemplates.com
%
% Original author:
% Frits Wenneker (http://www.howtotex.com) with extensive modifications by
% Vel (vel@LaTeXTemplates.com)
%
% License:
% CC BY-NC-SA 3.0 (http://creativecommons.org/licenses/by-nc-sa/3.0/)
%
%%%%%%%%%%%%%%%%%%%%%%%%%%%%%%%%%%%%%%%%%

%----------------------------------------------------------------------------------------
%	PACKAGES AND OTHER DOCUMENT CONFIGURATIONS
%----------------------------------------------------------------------------------------

\documentclass[twoside,twocolumn]{article}

\usepackage{blindtext} % Package to generate dummy text throughout this template 

\usepackage[sc]{mathpazo} % Use the Palatino font
\usepackage[T1]{fontenc} % Use 8-bit encoding that has 256 glyphs
\linespread{1.05} % Line spacing - Palatino needs more space between lines
\usepackage{microtype} % Slightly tweak font spacing for aesthetics

\usepackage[english]{babel} % Language hyphenation and typographical rules

\usepackage[hmarginratio=1:1,top=32mm,columnsep=20pt]{geometry} % Document margins
\usepackage[hang, small,labelfont=bf,up,textfont=it,up]{caption} % Custom captions under/above floats in tables or figures
\usepackage{booktabs} % Horizontal rules in tables

\usepackage{lettrine} % The lettrine is the first enlarged letter at the beginning of the text

\usepackage{enumitem} % Customized lists
\setlist[itemize]{noitemsep} % Make itemize lists more compact

\usepackage{abstract} % Allows abstract customization
\renewcommand{\abstractnamefont}{\normalfont\bfseries} % Set the "Abstract" text to bold
\renewcommand{\abstracttextfont}{\normalfont\small\itshape} % Set the abstract itself to small italic text

\usepackage{titlesec} % Allows customization of titles

\titleformat{\section}[block]{\large\scshape\centering}{\thesection.}{1em}{} % Change the look of the section titles
\titleformat{\subsection}[block]{\large}{\thesubsection.}{1em}{} % Change the look of the section titles

\usepackage{fancyhdr} % Headers and footers
\usepackage{graphicx}
\graphicspath{{pictures/}}
\DeclareGraphicsExtensions{.pdf,.png,.jpg}
\pagestyle{fancy} % All pages have headers and footers
\fancyhead{} % Blank out the default header
\fancyfoot{} % Blank out the default footer
\fancyfoot[RO,LE]{\thepage} % Custom footer text

\usepackage{titling} % Customizing the title section

\usepackage{hyperref} % For hyperlinks in the PDF

\usepackage {indentfirst} % For red string at first paragraph

%----------------------------------------------------------------------------------------
%	TITLE SECTION
%----------------------------------------------------------------------------------------

\setlength{\droptitle}{-4\baselineskip} % Move the title up

\pretitle{\begin{center}\Huge\bfseries} % Article title formatting
\posttitle{\end{center}} % Article title closing formatting

\begin{titlepage}
    \centering
    \textsc{Санкт-Петербургский политехнический университет Петра Великого}\\[3mm]
    \textsc{Институт компьютерных наук и технологий}\\[3mm]
    \textsc{Кафедра компьютерных систем и программных технологий}
	
	\vfill
	
	\textbf{Отчёт по лабораторной работе }\\[3mm]
	по курсу «Проектирование ОС и компонентов»\\[3mm]
	по теме «Драйвер символьного устройства»\\[41mm]
	
    \begin{flushright}
	\begin{minipage}{.35\textwidth}
		Выполнил студент гр. 13541/2:\\
		Волкова М. Д.\\[3mm]
		Проверил преподаватель:\\
		Душутина Е. В.
	\end{minipage}
    \end{flushright}
	
	\vfill

	Санкт-Петербург\\
	\the\year\ г.
\end{titlepage}


\renewcommand{\maketitlehookd}{
\begin{abstract}

The document contains an overview of the deep learning frameworks, the basic models and architectures of neural networks, and already existing deep learning laboratory works.

The document considers the tasks, models of deep learning and theoretical part. An analysis of tools for working with deep learning was conducted. The development was conducted in Python in a Jupyter Notebook.
\end{abstract}


}

\begin{document}
\maketitle

\section{Introduction.}

\par Currently, deep learning is an open-ended research area. With their success, deep learning is required to constantly increase the ability to go around data and relatively cheap graphics processors, allowing to build an effective calculation procedure.
In-depth training through sequential non-linear transformations, which, as a rule, are represented in the form of artificial neural networks. Today, they are used to solve such problems as prediction, pattern recognition and a number of others.
\par Deep learning is a subset of machine learning methods that use artificial neuron networks, built on the basis of an analogy with the structure of the neurons of the human brain. The term “deep” implies a large number of layers in the neural network.
\par Google, Microsoft, Amazon, Apple, Facebook and many other deep learning methods are used to analyze large data sets. Now this knowledge and skills went beyond purely academic research and became the property of large industrial companies.
\par The relevance of the topic of deep learning is confirmed by the regular appearance of articles on this topic.

\par The aim of the thesis is the development of teaching tools for the study of models of deep learning. In accordance with the purpose of the study, the following tasks were set:
\begin{itemize}
    \item analyze the main laboratory work;
    \item compare and choose technology for modeling deep neyron
networks;
    \item develop and test several laboratory works.
    \end{itemize}
\section{Review of existing deep learning technologies.}

With more and more businesses looking to scale up their operations, it has become integral for them to imbibe both machine learning as well as predictive analytics. AI coupled with the right deep learning framework has truly amplified the overall scale of what businesses can achieve and obtain within their domains.

The machine learning paradigm is continuously evolving. The key is to shift towards developing machine learning models that run on mobile in order to make applications smarter and far more intelligent. Deep learning is what makes solving complex problems possible. As put in this article, Deep Learning is basically Machine Learning on steroids. There are multiple layers to process features, and generally, each layer extracts some piece of valuable information.

Each framework is built in a different manner for different purposes. Here, we look at the 8 deep learning frameworks to give you a better idea of which framework will be the perfect fit or come handy in solving your business challenges.

\subsection{TensorFlow}
TensorFlow is arguably one of the best deep learning frameworks and has been adopted by several giants such as Airbus, Twitter, IBM, and others mainly due to its highly flexible system architecture.

The most well-known use case of TensorFlow has got to be Google Translate coupled with capabilities such as natural language processing, text classification/summarization, speech/image/handwriting recognition, forecasting, and tagging.

TensorFlow is available on both desktop and mobile and also supports languages such as Python, C++, and R to create deep learning models along with wrapper libraries.

TensorFlow comes with two tools that are widely used:

\begin{enumerate}
    \item TensorBoard for the effective data visualization of network modeling and performance.
    \item TensorFlow Serving for the rapid deployment of new algorithms/experiments while retaining the same server architecture and APIs. It also provides integration with other TensorFlow models, which is different from conventional practices and can be extended to serve other model and data types.
\end{enumerate}


If you're taking your first steps toward deep learning, it is a no-brainer to opt for TensorFlow given that is Python-based, is supported by Google, and comes loaded with documentation and walkthroughs to guide you.

\subsection{Caffe}
Caffe is a deep learning framework that is supported with interfaces like C, C++, Python, and MATLAB as well as the command line interface. It is well known for its speed and transposability and its applicability in modeling convolution neural networks (CNN). The biggest benefit of using Caffe's C++ library (comes with a Python interface) is the ability to access available networks from the deep net repository Caffe Model Zoo that are pre-trained and can be used immediately. When it comes to modeling CNNs or solving image processing issues, this should be your go-to library.

Caffe's biggest USP is speed. It can process over 60 million images on a daily basis with a single Nvidia K40 GPU. That's 1 ms/image for inference and 4 ms/image for learning-and more recent library versions are faster still.

Caffe is a popular deep learning network for visual recognition. However, Caffe does not support fine-granular network layers like those found in TensorFlow or CNTK. Given the architecture, the overall support for recurrent networks, and language modeling it's quite poor, and establishing complex layer types has to be done in a low-level language.


\subsection{Microsoft Cognitive Toolkit/CNTK}

Popularly known for easy training and the combination of popular model types across servers, the Microsoft Cognitive Toolkit (previously known as CNTK) is an open-source deep learning framework to train deep learning models. It performs efficient convolution neural networks and training for image, speech, and text-based data. Similar to Caffe, it is supported by interfaces such as Python, C++, and the command line interface.

Given its coherent use of resources, the implementation of reinforcement learning models or generative adversarial networks (GANs) can be done easily using the toolkit. It is known to provide higher performance and scalability as compared to toolkits like Theano or TensorFlow while operating on multiple machines.

Compared to Caffe, when it comes to inventing new complex layer types, users don't need to implement them in a low-level language due to the fine granularity of the building blocks. The Microsoft Cognitive Toolkit supports both RNN and CNN types of neural models and thus is capable of handling images, handwriting, and speech recognition problems. Currently, due to the lack of support on ARM architecture, its capabilities on mobile are fairly limited.


\subsection{Torch/PyTorch}
Torch is a scientific computing framework that offers wide support for machine learning algorithms. It is a Lua-based deep learning framework and is used widely amongst industry giants such as Facebook, Twitter, and Google. It employs CUDA along with C/C++ libraries for processing and was basically made to scale the production of building models and provide overall flexibility.

As of late, PyTorch has seen a high level of adoption within the deep learning framework community and is considered to be a competitor to TensorFlow. PyTorch is basically a port to the Torch deep learning framework used for constructing deep neural networks and executing tensor computations that are high in terms of complexity.

As opposed to Torch, PyTorch runs on Python, which means that anyone with a basic understanding of Python can get started on building their own deep learning models.

Given PyTorch framework's architectural style, the entire deep modeling process is far simpler as well as transparent compared to Torch.

\subsection{Keras}

Well known for being minimalistic, the Keras neural network library (with a supporting interface of Python) supports both convolutional and recurrent networks that are capable of running on either TensorFlow or Theano. The library is written in Python and was developed keeping quick experimentation as its USP.

Due to the fact that the TensorFlow interface is a tad bit challenging coupled with the fact that it is a low-level library that can be intricate for new users, Keras was built to provide a simplistic interface for the purpose of quick prototyping by constructing effective neural networks that can work with TensorFlow.

Lightweight, easy to use, and really straightforward when it comes to building a deep learning model by stacking multiple layers: that is Keras in a nutshell. These are the very reasons why Keras is a part of TensorFlow's core API.

The primary usage of Keras is in classification, text generation and summarization, tagging, and translation, along with speech recognition and more. If you happen to be a developer with some experience in Python and wish to dive into deep learning, Keras is something you should definitely check out.


\subsection{Conclusion}

It is evident that the advent of deep learning has initiated many practical use cases of machine learning and artificial intelligence. Breaking down tasks in the simplest of ways in order to assist machines in the most efficient manner has been made possible by deep learning.

That being said, which deep learning framework from the above list would best suit your business requirements? The answer to that lies on a number of factors, however, if you are looking to just get started, then a Python-based deep learning framework like TensorFlow or Chainer is ideal.

If you're looking for something more, in that case - speed, resource requirement, and usage along with the coherence of the trained model should always be considered prior to picking out the best deep learning framework for your business needs.

At Maruti Techlabs, we extensively use TensorFlow and Keras for our client requirements - one of them being processing images for an online car marketplace. Images are recognized, identified and differentiated while also understanding the objects in the image. The algorithm was put in place mainly to assess and flag images that were not associated with cars and thus, maintaining the quality and precision of image related data was taken care of.

There you have it! 8 of the best deep learning frameworks for data science enthusiasts. We hope you liked this post and learned something new and useful. If you did, feel free to give it some claps.
\section{Neural network models}

If you've heard about Artificial Intelligence, Machine Learning, or Deep Learning recently, then you might have heard of a Neural Network.

Neural Networks are a key piece of some of the most successful machine learning algorithms. The development of neural networks have been key to teaching computers to think and understand the world in the way that humans do. Essentially, a neural network emulates the human brain. Brains cells, or neurons, are connected via synapses. This is abstracted as a graph of nodes (neurons) connected by weighted edges (synapses).

So let's dive in. What is a neural network? The human brain consists of 100 billion cells called neurons, connected together by synapses. If sufficient synaptic inputs fire to a neuron, that neuron will also fire. We call this process "thinking". We can model this process by creating a neural network on a computer. A neural network has input and output neurons, which are connected by weighted synapses. The weights affect how much of the forward propagation goes through the neural network. The weights can then be changed during the back propagation-this is the part where the neural network is now learning. This process of forward propagation and backward propagation is conducted iteratively on every piece of data in a training data set. The greater the size of the data set and the greater the variety of data set that there is, the more that the neural network will learn, and the better that the neural network will get at predicting outputs.


\subsection{Perceptrons}
 
Considered the first generation of neural networks, Perceptrons are simply computational models of a single neuron. Perceptron was originally coined by Frank Rosenblatt in his paper, "The perceptron: a probabilistic model for information storage and organization in the brain" (1956). Also called feed-forward neural network, perceptron feeds information from the front to the back. Training perceptrons usually requires back-propagation, giving the network paired datasets of inputs and outputs. Inputs are sent into the neuron, processed, and result in an output. The error being back propagated is often some variation of the difference between the input and the output. Given that the network has enough hidden neurons, it can theoretically always model the relationship between the input and output. Practically their use is a lot more limited but they are popularly combined with other networks to form new networks.


However, Perceptrons do have limitations: If you choose features by hand and you have enough features, you can do almost anything. For binary input vectors, we can have a separate feature unit for each of the exponentially many binary vectors and so we can make any possible discrimination on binary input vectors. But once the hand-coded features have been determined, there are very strong limitations on what a perceptron can learn.
 
\subsection{Convolutional Neural Networks}

In 1998, Yann LeCun and his collaborators developed a really good recognizer for handwritten digits called LeNet. It used back propagation in a feedforward net with many hidden layers, many maps of replicated units in each layer, pooling of the outputs of nearby replicated units, a wide net that can cope with several characters at once even if they overlap, and a clever way of training a complete system, not just a recognizer. Later it is formalized under the name convolutional neural networks (CNNs). Fun fact: This net was used for reading ~10\% of the checks in North America.


Convolutional Neural Networks are quite different from most other networks. They are primarily used for image processing but can also be used for other types of input such as as audio. A typical use case for CNNs is where you feed the network images and the network classifies the data. CNNs tend to start with an input "scanner" which is not intended to parse all the training data at once. For example, to input an image of 100 x 100 pixels, you wouldn't want a layer with 10 000 nodes. Rather, you create a scanning input layer of say 10 x 10 which you feed the first 10 x 10 pixels of the image. Once you passed that input, you feed it the next 10 x 10 pixels by moving the scanner one pixel to the right.


This input data is then fed through convolutional layers instead of normal layers, where not all nodes are connected to all nodes. Each node only concerns itself with close neighboring cells. These convolutional layers also tend to shrink as they become deeper, mostly by easily divisible factors of the input. Besides these convolutional layers, they also often feature pooling layers. Pooling is a way to filter out details: a commonly found pooling technique is max pooling, where we take say 2 x 2 pixels and pass on the pixel with the most amount of red. If you want to dig deeper into CNNs, read Yann LeCun's original paper - "Gradient-based learning applied to document recognition" (1998).

 \subsection{Recurrent Neural Network}

To understand RNNs, we need to have a brief overview on sequence modeling. When applying machine learning to sequences, we often want to turn an input sequence into an output sequence that lives in a different domain; for example, turn a sequence of sound pressures into a sequence of word identities. When there is no separate target sequence, we can get a teaching signal by trying to predict the next term in the input sequence. The target output sequence is the input sequence with an advance of 1 step. This seems much more natural than trying to predict one pixel in an image from the other pixels, or one patch of an image from the rest of the image. Predicting the next term in a sequence blurs the distinction between supervised and unsupervised learning. It uses methods designed for supervised learning, but it doesn't require a separate teaching signal.

One big problem with RNNs is the vanishing (or exploding) gradient problem where, depending on the activation functions used, information rapidly gets lost over time. Intuitively this wouldn't be much of a problem because these are just weights and not neuron states, but the weights through time is actually where the information from the past is stored; if the weight reaches a value of 0 or 1 000 000, the previous state won't be very informative. RNNs can in principle be used in many fields as most forms of data that don't actually have a timeline (i.e. unlike sound or video) can be represented as a sequence. A picture or a string of text can be fed one pixel or character at a time, so the time dependent weights are used for what came before in the sequence, not actually from what happened x seconds before. In general, recurrent networks are a good choice for advancing or completing information, such as autocompletion.


\subsection{Long/Short Term Memory Network}
Hochreiter & Schmidhuber (1997) solved the problem of getting a RNN to remember things for a long time by building what known as long-short term memory networks (LSTMs). LSTMs networks try to combat the vanishing / exploding gradient problem by introducing gates and an explicitly defined memory cell. The memory cell stores the previous values and holds onto it unless a "forget gate" tells the cell to forget those values. LSTMs also have a "input gate" which adds new stuff to the cell and an "output gate" which decides when to pass along the vectors from the cell to the next hidden state.

Recall that with all RNNs, the values coming in from X train and H previous are used to determine what happens in the current hidden state. And the results of the current hidden state (H current) are used to determine what happens in the next hidden state. LSTMs simply add a cell layer to make sure the transfer of hidden state information from one iteration to the next is reasonably high. Put another way, we want to remember stuff from previous iterations for as long as needed, and the cells in LSTMs allow this to happen. LSTMs have been shown to be able to learn complex sequences, such as writing like Shakespeare or composing primitive music.
\section{Discussion.}

The purpose of this work was to introduce the basics and concepts of deep learning. The biggest problem that can be encountered in the study of deep learning is the cost of equipment.


Since this project is educational, it was decided to study more architecture problems, but rather superficially. In the future, in order to make the course more complete, it was possible to add more theory and increase the number of problems to be solved.


As a result, we have a working system ready for use.
\section{Results}

The course is designed to be performed with any operating system. To work with textbooks, you need to install Python 3 and Jupyter. The number of third-party libraries is minimized.

The notebooks were designed so that, using their example, students could learn how to extract features from heterogeneous data, what problems there are and how to solve them. They will also learn how to reduce the tasks to the formal statements of the machine learning problem and understand how to check the quality of the constructed model on different data.

Having completed this course, a student can gain knowledge about common types of applied problems and be able to find schemes for solving them.

The course contains three laboratory works that solve problems: classification, localization of images and time series prediction.
\section{Conclusion.}

As a result of the bachelor's work, a study was made of deep learning technologies: the existing and new technologies that were applied, as well as various deep learning programs were reviewed. The analysis of laboratory work laid out in open access was made.
TensorFlow and Keras libraries were chosen among the many software tools for deep learning. They meet the needs of the course being developed, namely, it is suitable for solving assigned tasks.
To visualize the search results, Jupyter Notebook graphic web shell was selected. The design was carried out in a graphic web shell Jupyter Notebook. In it were the laboratory works themselves.

The practical significance of the development of an educational-methodical complex of laboratory works lies in the possibility of introducing it into the educational process.
Performing a course of laboratory work on the subject “study of models of deep learning” consists of 3 works. Each laboratory focuses on the study of various tasks: classical fictions, image localization and time series prediction. Upon their completion, students will learn to reduce various tasks to formal statements of machine learning problems and to apply the necessary technologies and schemes for their solutions.
The tasks of the bachelor's work were fulfilled, and the goals were achieved.

\begin{thebibliography}

\begin{itemize}
    \item [1.] McCulloch W.S. Pitts W. A Logical Calculus of the Ideas Immanent in Nervous Activity. - The Bulletin of Mathematical Biophysics, 1943. - 133 p.
    
    \item [2.] NVIDIAElectronic resource. - URL: https://www.nvidia.ru/deep-learning-ai/education/.

    \item [3.] Ballard, D. H., Hinton, G. E., and Sejnowski, T. J. (1983). Parallel vision computation.Nature

    \item [4.] LeCun, Y., Bottou, L., Orr, G. B., and Muller, K.-R. (1998a). Efficient backprop. InNeural Networks, Tricks of the Trade, Lecture Notes in Computer Science LNCS 1524.Springer Verlag.

    \item [5.] Roweis, S., Saul, L., and Hinton, G. (2002). Global coordination of local linear models. InT. Dietterich, S. Becker, and Z. Ghahramani, editors, Advances in Neural InformationProcessing Systems 14 (NIPS'01), Cambridge, MA. MIT Press.

\end{itemize}

\end{thebibliography}

\end{document}

\addcontentsline{toc}{chapter}{Заключение}
\chapter*{Заключение}

Cisco Packet Tracer \cite{cite-cis} представляет собой удобный инструмент для построения макетов сетей, предоставляя множество инструментов для настройки узлов сети. Большим преимуществом можно считать множество готовых компонентов, которые особенно удобны для реализации сетевых сервисов.

Данную работу можно было выполнить в другой системе виртуализации -- VMware Workstation \cite{cite-vm}, однако, она не является специализированной именно для моделирования компьютерных сетей, что накладывает некоторые ограничения, а также вынуждает в ручной установке операционных систем и настройке сетевых сервисов. Кроме того, специализированные программы по типу Cisco Packet Tracer \cite{cite-cis} оценивают задержку сигнала, в зависимости от расстояния и множества других паразитных параметров. 

%------------------------------------------------------------------------------

%\addcontentsline{toc}{section}{Список литературы}
%\bibliography{thesis}
%\bibliographystyle{ugost2008}

%\nocite{phill}
%\nocite{Ren}

\bibliography{thesis}
\bibliographystyle{ugost2008}

\end{document}
